% !TeX root = ../教材题库_全部条目_转写工作稿.tex
% Source: XJ-B1-C04-HW-P125-Y0545
% 第4章 / 4.3 对数函数 / 教材 p.125 / PDF 第132页
% OCR未校对：公式与文字均可能有误；未完成数学TeX转写；禁止直接发给学生
% 原题及相邻原解截图（核对用）：../题目原图/XJ-B1-C04-HW-P125-Y0545-01.webp
% 本文件只在转写完成后作为题干使用；OCR 条目尚不是可靠题干。
\begin{bankitem}{习题4.3 第7题}{第4章 / 4.3 对数函数 / 教材 p.125 / PDF 第132页}{XJ-B1-C04-HW-P125-Y0545}{OCR未校对：公式与文字均可能有误；未完成数学TeX转写；禁止直接发给学生}
\begin{ocrdraft}
[教材p.125 / PDF第132页]
\par
7．溶液的酸碱度（pH）是由公式pH=—lgCH]计量的，其中[H]表示溶液中
\par
氢离子的浓度，单位为mol/L.pH的范围为0、14．
\par
(I)CH]为0，1mol/L，0．01mol/L的溶液的pH分别是多少？
\par
（2）随着氢离子浓度的降低，pH会发生怎样的变化？
\par
（3）试确定橙汁（pH一3．5）的氢离子浓度．
\par
（4）试确定人体血液在pH=7.4时的氢离子浓度．
\end{ocrdraft}
\end{bankitem}
